\documentclass[a4paper,10pt]{article}

% ---------- layout and fonts ----------
\usepackage[a4paper,left=1.7cm,right=1.7cm,top=1.3cm,bottom=1.2cm]{geometry}
\usepackage[T1]{fontenc}
\usepackage[sfdefault]{plex-sans}
\usepackage[spanish,es-noshorthands,es-nodecimaldot]{babel}
\usepackage{xcolor}
\usepackage{enumitem}
\usepackage[hidelinks]{hyperref}
\input{glyphtounicode}
\pdfgentounicode=1

\definecolor{ink}{HTML}{1B1B1B}
\definecolor{muted}{HTML}{5F5F5F}
\definecolor{accent}{HTML}{1F4E8C}
\definecolor{rule}{HTML}{D9D4CB}

\pagestyle{empty}
\setlength{\parindent}{0pt}
\linespread{1.08}
\raggedright
\hyphenpenalty=10000
\exhyphenpenalty=10000
\color{ink}

\hypersetup{pdftitle={Javier Herrero Pérez, currículum}, pdfauthor={Javier Herrero Pérez}}

% ---------- building blocks ----------
\newcommand{\sectionhead}[1]{%
  \vspace{8pt}{\color{accent}\fontsize{11.5}{14}\selectfont\bfseries #1}\par
  \vspace{3pt}{\color{rule}\hrule height 0.5pt}\vspace{6pt}}

% role, dates, organisation, place
\newcommand{\entry}[4]{%
  {\bfseries #1}\hfill{\color{muted}#2}\par
  {\color{muted}#3}\hfill{\color{muted}#4}\par\vspace{1pt}}

% title, link, date and context, description, tools
\newcommand{\project}[5]{%
  {\bfseries\href{#2}{#1}}\hfill{\color{muted}#3}\par
  #4\par
  {\color{muted}\small #5}\par\vspace{5pt}}

% label, items
\newcommand{\skill}[2]{\makebox[3.7cm][l]{\color{muted}#1}#2\par\vspace{1pt}}

\setlist[itemize]{leftmargin=11pt, itemsep=1pt, topsep=2pt, parsep=0pt,
  label={\color{muted}\raisebox{0.15ex}{\scriptsize\textbullet}}}

\begin{document}

% ---------- cabecera ----------
{\fontsize{22}{26}\selectfont\bfseries Javier Herrero Pérez}\par
\vspace{3pt}
{\large\color{muted} Científico de datos}\par
\vspace{6pt}
{\small
Salamanca, España\quad
656\,181\,225\quad
\href{mailto:jherrerop88@gmail.com}{jherrerop88@gmail.com}\quad
\href{https://www.linkedin.com/in/javier-herrero-p\%C3\%A9rez-ba672b271/}{\color{accent}LinkedIn}\quad
\href{https://github.com/Jheivy}{\color{accent}github.com/Jheivy}\quad
\href{https://jheivy.github.io/?lang=es}{\color{accent}jheivy.github.io}\par}

% ---------- perfil ----------
\sectionhead{Perfil}
Científico de datos con formación en Física. Trabajo en Whizdom.ai y estoy terminando el Máster en Ciencia de Datos de la Universidad de Valencia, a falta del TFM. Me interesan el machine learning en producción, la inferencia causal y los modelos de lenguaje.

% ---------- experiencia ----------
\sectionhead{Experiencia}
\entry{AI Researcher \& Developer}{ago 2026 -- actualidad}{Whizdom.ai}{En remoto}
\begin{itemize}
  \item Investigación y desarrollo de modelos de machine learning para un producto de analítica de iGaming, desde la preparación de los datos hasta producción.
  \item Datos a gran escala y MLOps reproducible: Python, SQL, seguimiento de experimentos, contenedores.
\end{itemize}
\vspace{6pt}
\entry{Prácticas en ciencia de datos}{jul -- ago 2024}{AIR Institute}{España}
\begin{itemize}
  \item Proyectos de machine learning aplicado sobre datos de sensores ambientales y transcripciones de entrevistas: clasificación, NLP y explicabilidad de modelos.
\end{itemize}

% ---------- formación ----------
\sectionhead{Formación}
\entry{Máster en Ciencia de Datos}{sept 2025 -- actualidad}{Universidad de Valencia, nota media 8,86/10, TFM en curso}{España}
\vspace{4pt}
\entry{Grado en Física}{sept 2021 -- jul 2025}{Universidad de Salamanca, nota media 7,25/10}{España}

% ---------- proyectos ----------
\sectionhead{Proyectos destacados}
\project{Clickbait en la prensa española}{https://github.com/hepejav/ProyectoNLP}{may 2026, equipo de 7}
  {Comparamos titulares con el texto de la noticia en nueve medios españoles con embeddings y modelos de temas, y construimos un agente con LangChain que extrae y etiqueta titulares nuevos.}
  {Python, Transformers, LangChain, UMAP}
\project{CalorIA: calorías a partir de una foto}{https://github.com/hepejav/CalorIA}{mar 2026, equipo de 2}
  {Un modelo de visión y lenguaje (Qwen2-VL) identifica alimentos y cantidades en la foto de un plato; los nutrientes salen de la base de datos del USDA y se comparan con el objetivo diario.}
  {Python, Hugging Face, LLMs}
\project{¿Sube la formación laboral los ingresos?}{https://github.com/hepejav/An-lisis-de-los-datos-de-Lalonde-1986-y-Dehejia-y-Wahba-1999-}{feb 2026}
  {Análisis causal de los datos de Lalonde con propensity score, trimming y regresión ajustada. Efecto estimado de unos +2.200 USD en los ingresos de 1978.}
  {R, inferencia causal}
\project{Abandono de clientes de telecomunicaciones}{https://github.com/hepejav/Telco_Customer_Chrun}{ene 2026}
  {Clasificador de abandono con variables de anomalías de facturación y un ensemble por votación con umbral ajustado. AUC de 0,85.}
  {Python, XGBoost, LightGBM, scikit-learn}
\project{Optimización de carteras con AutoARIMA}{https://github.com/Jheivy/Portfolio-Optimization-ARIMA-LogUtility}{ene 2026}
  {Predicciones a un paso sin sesgo de anticipación para carteras de media-varianza y utilidad logarítmica. Máxima nota.}
  {R, forecast, quadprog}

% ---------- herramientas ----------
\sectionhead{Herramientas}
\skill{Programación}{Python, SQL, R}
\skill{Modelado}{scikit-learn, LightGBM, XGBoost, SHAP, análisis de supervivencia, inferencia causal}
\skill{Deep learning, NLP}{PyTorch, TensorFlow, Hugging Face, LangChain}
\skill{Datos y MLOps}{DuckDB, AWS S3, MLflow, Docker, FastAPI, Git}
\skill{Idiomas}{Español (nativo), inglés (B1+, nivel profesional)}

\end{document}
